% from ega4-2-fr.tex
\oldpage[IV-2]{187}

\begin{remarks}[7.1.12]
\label{IV.2.s7.1.12-fr}
\begin{enumerate}
\item[(i)] On ignore si, réciproquement, un anneau local noethérien
universellement caténaire est toujours formellement caténaire.
\item[(ii)] Un anneau local noethérien $A$ \emph{de dimension 1} est nécessairement
équidimensionnel, son idéal maximal ne pouvant être minimal (car dans ce
cas $A$ serait de dimension 0). Comme $\dim(\widehat A)=1$, cela montre que
$A$ est même \emph{formellement équidimensionnel}, et \emph{a fortiori}
formellement caténaire (donc \emph{universellement caténaire}). Par contre,
l'anneau local de dimension 2 défini dans (5.6.11), qui est caténaire mais
non universellement caténaire, est \emph{a fortiori} non formellement
caténaire.
\end{enumerate}
\end{remarks}

\begin{corollary}[7.1.13]
\label{IV.2.s7.1.13-fr}
Soient $Y$ un préschéma localement noethérien irréductible de dimension 1,
$X$ un préschéma irréductible, $f:X\to Y$ un morphisme dominant de type
fini, $\xi$, $\eta$ les points génériques de $X$ et $Y$ respectivement.
Alors, pour tout $y\in f(X)$, les dimensions des composantes irréductibles
de $f^{-1}(y)$ sont toutes égales à
$\deg.\operatorname{tr}_{k(\eta)}k(\xi)$.
\end{corollary}

\begin{proof}
Par hypothèse, $f^{-1}(\eta)$ est irréductible de point générique $\xi$ et
de dimension $\deg.\operatorname{tr}_{k(\eta)}k(\xi)$ (5.2.1). Comme $Y$
est irréductible et de dimension 1, on a $\dim(\mathscr O_y)=1$ pour tout
$y\neq\eta$, et $\mathscr O_y$ est universellement caténaire en vertu de
(7.1.12, (ii)). Si $y\in f(X)$ et si $z$ est un point générique d'une
composante irréductible $Z$ de $f^{-1}(y)$, on a donc, par (5.6.5)
\[
\dim(Z)=1+\dim(f^{-1}(\eta))-\dim(\mathscr O_{X,z}).
\]
Mais on a $\dim(\mathscr O_{X,z})\leqslant\dim(\mathscr O_y)$ (0,
16.3.9), et d'autre part, comme $z$ n'est pas le point générique de $X$,
$\dim(\mathscr O_{X,z})>0$; donc $\dim(\mathscr O_{X,z})=1$ et
$\dim(Z)=\dim(f^{-1}(\eta))$.
\end{proof}

\subsection{Anneaux strictement formellement caténaires}
\label{subsection:IV.7.2-fr}

\begin{env}[7.2.1]
\label{IV.2.s7.2.1-fr}
\emph{Notations.} --- Pour un anneau intègre noethérien $A$, nous noterons
$A^{(1)}$ l'intersection des anneaux locaux $A_{\mathfrak p}$, où
$\mathfrak p$ parcourt l'ensemble des idéaux premiers de $A$ de hauteur 1
(5.10.17). Si $A$ est un anneau local noethérien intègre, nous noterons
$A^{(\omega)}$ l'intersection des $A_{\mathfrak p}$ où cette fois
$\mathfrak p$ parcourt l'ensemble de tous les idéaux premiers de $A$
distincts de l'idéal maximal.

Nous dirons qu'un anneau local noethérien $A$ est \emph{strictement
équidimensionnel} si, pour tout idéal premier
$\mathfrak p\in\operatorname{Ass}(A)$, on a
$\dim(A/\mathfrak p)=\dim(A)$; il revient au même de dire que $A$ est
équidimensionnel et \emph{sans idéaux premiers associés immergés}.
\end{env}

\begin{example}[7.2.1.1]
\label{IV.2.s7.2.1.1-fr}
Soit $A$ un anneau local noethérien \emph{de dimension 1}. Alors les
conditions suivantes sont équivalentes:
\begin{enumerate}
\item[a)] $A$ est sans idéaux premiers associés immergés.
\item[$a'$)] $\widehat A$ est sans idéaux premiers associés immergés.
\item[b)] $A$ est strictement équidimensionnel.
\item[$b'$)] $\widehat A$ est strictement équidimensionnel.
\item[c)] $A$ est un anneau de Cohen--Macaulay.
\end{enumerate}
\end{example}

En effet a) signifie que l'idéal maximal $\mathfrak m$ de $A$ n'est pas
associé à $A$, donc les idéaux premiers associés à $A$ sont les idéaux
minimaux de $A$, tous distincts de $\mathfrak m$, et

